\documentclass[conference]{IEEEtran}

\IEEEoverridecommandlockouts

\usepackage{cite}
\usepackage{amsmath,amssymb,amsfonts}
\usepackage{graphicx}
\usepackage{textcomp}
\usepackage{xcolor}
\usepackage{booktabs}
\usepackage{url}
\usepackage{array}
\usepackage{placeins}

\graphicspath{{./}{../}}

\def\BibTeX{{\rm B\kern-.05em{\sc i\kern-.025em b}\kern-.08em
T\kern-.1667em\lower.7ex\hbox{E}\kern-.125emX}}

\newcommand{\figplaceholder}[1]{%
\fbox{\parbox[c][1.55in][c]{0.92\columnwidth}{\centering #1}}%
}
\newcommand{\figimage}[2]{%
\IfFileExists{#2}{\includegraphics[width=#1]{#2}}{%
\IfFileExists{../#2}{\includegraphics[width=#1]{../#2}}{%
\figplaceholder{Place #2 here.}}}%
}

\begin{document}

\title{GeoRisk-IQ: A Hybrid Spatial-Financial Framework for Bank Branch Opportunity Scoring and Failure-Risk Validation}

\author{
\IEEEauthorblockN{Harapriya Bisoyee}
\IEEEauthorblockA{
Tech Mahindra Americas\\
harapriyab2022@gmail.com}
}

\maketitle

\begin{abstract}
We present GeoRisk-IQ, a geospatial machine-learning framework for bank
branch opportunity scoring and bank-failure risk validation. The framework
maps Federal Deposit Insurance Corporation (FDIC) institution records to an
H3-style spatial grid, derives local and neighborhood density features, and
combines these features with cell-level financial indicators. GeoRisk-IQ
supports two related tasks: ranking underserved opportunity zones and
validating spatial-financial signals against historical bank-failure labels.
Using 27,832 FDIC institution records aggregated into 6,086 spatial cells,
the opportunity scorer obtains mean absolute error of 0.0081 and
$R^2=0.9991$ against a spatial opportunity proxy. For failure-risk
validation, a hybrid spatial-financial gradient boosting model achieves
10-fold cross-validation ROC-AUC of 0.8341 with a bootstrap 95\% confidence
interval of [0.8226, 0.8454]. A permutation test shows the hybrid model
improves over the spatial-only baseline ($\Delta$AUC = +0.0062,
$p=0.0013$) and substantially improves over the financial-only baseline
($\Delta$AUC = +0.2312, $p<0.0001$). These results indicate that spatial
neighborhood structure is a useful and interpretable signal for financial
location intelligence.
\end{abstract}

\begin{IEEEkeywords}
geospatial AI, banking analytics, spatial risk, opportunity scoring,
FDIC, gradient boosting, H3-style spatial indexing
\end{IEEEkeywords}

\section{Introduction}
Branch network decisions require a joint view of geography, local market
supply, and institutional financial condition. A cell with few branches may
be an attractive opportunity if it is adjacent to dense banking corridors;
conversely, a branch operating in an isolated or saturated local market may
exhibit different risk dynamics than a branch with similar financial ratios
in a healthier surrounding market. Conventional bank-risk models typically
focus on institution-level balance-sheet variables and underuse the spatial
structure of local banking markets.

GeoRisk-IQ addresses this gap with a spatial-financial machine-learning
pipeline. The system converts point-level FDIC banking records into a
regular spatial grid, engineers neighborhood features over multiple
$k$-ring radii, and trains models for opportunity scoring and failure-risk
validation. The implementation is exposed through an interactive Streamlit
dashboard that supports opportunity-zone maps, branch-density maps, model
comparison, temporal validation, and feature-importance analysis.

The contributions of this work are:
\begin{itemize}
    \item A reproducible spatial feature pipeline for FDIC institution data
    using local density, neighboring density, spatial ratios, isolation, and
    financial aggregates.
    \item A branch opportunity scorer that ranks cells with low local branch
    supply and high surrounding banking density.
    \item A four-way model comparison for bank-failure validation:
    logistic regression, spatial-only gradient boosting, financial-only
    gradient boosting, and hybrid spatial-financial gradient boosting.
    \item A validation protocol with 10-fold stratified cross-validation,
    temporal holdout, bootstrap confidence intervals, permutation tests,
    ablation, and SHAP-based feature attribution.
\end{itemize}

\section{Related Work}
Bank failure prediction has traditionally relied on CAMELS-style financial
ratios and supervisory indicators, including capitalization, asset quality,
earnings, liquidity, and sensitivity to market risk. Logistic regression,
survival models, and tree ensembles have all been applied to bank distress
and bankruptcy prediction \cite{cole1998,betz2014,tian2015}. These methods
primarily model institutions as independent entities and often omit local
market geography.

Geospatial financial-risk studies have used administrative units such as ZIP
codes, counties, or census tracts to model mortgage default, lending access,
or local economic distress. Such units can introduce boundary artifacts and
scale dependence. Spatial indexing systems such as H3 \cite{h3} offer a
more regular representation for locality-preserving aggregation and
neighbor operations. GeoRisk-IQ follows this principle through an H3-style
grid abstraction and evaluates how local and surrounding density features
contribute to banking opportunity and risk modeling.

\subsection{Bank Distress and Failure Prediction}
Classical failure-prediction systems use accounting ratios, supervisory
ratings, and macroeconomic controls to estimate institutional distress.
Cole and Gunther \cite{cole1998} compare off-site monitoring systems with
on-site examination information, showing that financial ratios can identify
distress but require frequent updating. Later work extends these ideas with
survival analysis and nonlinear machine-learning models
\cite{betz2014,tian2015}. These approaches are strong baselines because
they capture institution-level health, but they do not directly represent
local market structure.

This omission matters for branch-level and market-cell analysis. Two
institutions with similar return-on-assets may operate in very different
geographic contexts: one may be isolated in a thinly served area, while
another may be embedded in a dense and competitive corridor. GeoRisk-IQ
therefore treats geography as a first-class modeling signal rather than a
post-hoc visualization layer.

\subsection{Spatial Aggregation for Financial Access}
Spatial aggregation is widely used in financial inclusion, mortgage
analytics, urban economics, and service-access studies. Administrative
units such as counties and ZIP codes are convenient but can create
scale-dependent conclusions because their boundaries are historical and
political rather than model-driven. Grid-based spatial indexing offers an
alternative: every observation is mapped into a repeatable spatial unit,
and neighboring units can be queried consistently. This supports both
visual analysis and machine-learning feature construction.

\subsection{Positioning}
GeoRisk-IQ differs from prior bank-risk work in three ways. First, it
combines spatial and financial signals in the same feature matrix. Second,
it evaluates spatial contribution through ablation and model comparison
rather than relying only on map visualization. Third, it supports an
operational dashboard that allows users to inspect the same cells used by
the model.

\begin{table}[t]
\caption{Positioning Relative to Related Modeling Approaches}
\label{tab:positioning}
\centering
\begin{tabular}{p{0.28\columnwidth}p{0.24\columnwidth}p{0.30\columnwidth}}
\toprule
Approach & Main signal & Limitation addressed by GeoRisk-IQ \\
\midrule
CAMELS-style models & Financial ratios & Add explicit spatial context \\
County/ZIP models & Administrative geography & Use repeatable grid cells \\
Branch-count maps & Visual density & Train and validate predictive models \\
GeoRisk-IQ & Spatial + financial & Joint scoring and validation \\
\bottomrule
\end{tabular}
\end{table}

\section{Data}
\subsection{FDIC Institution Registry}
The FDIC BankFind Suite provides institution records containing coordinates,
certificate identifiers, location metadata, assets, and office counts. The
current GeoRisk-IQ run loads 27,832 institutions from the local FDIC
institution file. Records without valid latitude and longitude are removed
before grid construction.

Institution coordinates are treated as the primary spatial signal. The
pipeline first normalizes field names, coerces latitude and longitude into
numeric values, and drops records that cannot be mapped. Asset fields are
coerced to numeric values when available. These preprocessing steps are
deliberately conservative because invalid coordinates can create false
cells and distort neighborhood density.

\subsection{FDIC Failure Records}
Historical failure records are loaded from \texttt{FDICFailure.txt}. When
certificate identifiers are available, failures are linked to institution
coordinates and then mapped to spatial cells. In the current validation
dataset, 6,086 cells are generated and 1,403 cells receive at least one
historical failure label.

\subsection{FDIC Financial Records}
Financial features are loaded from \texttt{FDICFinancialsFull.txt}. The
current reproducible experiment uses the latest available quarterly FDIC
financial snapshot and aggregates return on assets, net income, and assets
to the spatial-cell level. This design keeps the experiment lightweight and
reproducible, while leaving room for future work on historical lagged
financial time series.

\begin{table}[t]
\caption{Datasets Used in the Current Reproducible Run}
\label{tab:data_summary}
\centering
\begin{tabular}{lrl}
\toprule
Source file & Records & Role \\
\midrule
\texttt{Geospatiallocation.txt} & 27,832 & Institution locations \\
\texttt{FDICFailure.txt} & 4,114 & Failure events \\
\texttt{FDICFinancialsFull.txt} & 4,408 & Latest-quarter financials \\
\bottomrule
\end{tabular}
\end{table}

\subsection{Data Quality Considerations}
The main data-quality risk is geocoding reliability. Failure records do not
always contain direct branch-level coordinates, so GeoRisk-IQ links failure
events to institution records through certificate identifiers when possible.
Cells without reliable coordinate mappings are excluded from the labeled
validation grid. A second consideration is temporal alignment: the current
financial file is a latest-quarter snapshot, while failures span historical
periods. This makes the current financial-only model a weak historical
baseline and motivates future time-aware financial features.

\section{Feature Engineering}
\subsection{Spatial Cell Construction}
Each institution coordinate $(lat_i, lon_i)$ is assigned to a deterministic
spatial cell $c_i$. For each cell $c$, branch count is defined as
\begin{equation}
    B_c = \sum_i \mathbb{I}(c_i = c),
\end{equation}
and total assets are aggregated as
\begin{equation}
    A_c = \sum_i A_i \mathbb{I}(c_i = c).
\end{equation}

\subsection{Neighborhood Density}
For each cell $c$, GeoRisk-IQ computes average neighboring density across
multiple radii. Let $\mathcal{N}_k(c)$ denote the set of cells within
$k$ grid steps of $c$. The $k$-ring neighbor density is
\begin{equation}
    \bar{B}_{c,k} =
    \frac{1}{|\mathcal{N}_k(c)|}\sum_{j \in \mathcal{N}_k(c)} B_j.
\end{equation}
The current feature set includes $k \in \{1,2,3\}$, density ratios,
an isolation indicator, total asset features, and financial means.

\begin{table}[t]
\caption{Feature Groups Used by GeoRisk-IQ}
\label{tab:feature_groups}
\centering
\begin{tabular}{lll}
\toprule
Group & Examples & Purpose \\
\midrule
Local supply & branch\_count, asset\_total & Current cell intensity \\
Neighbor density & neighbor\_avg\_k1--k3 & Surrounding activity \\
Ratios & density\_ratio\_k1--k2 & Relative saturation \\
Isolation & is\_isolated & Sparse-market flag \\
Financial & ROA, net income, assets & Cell financial health \\
\bottomrule
\end{tabular}
\end{table}

\subsection{Opportunity Target}
The opportunity-scoring task uses a spatial proxy target:
\begin{equation}
    O_c = \frac{\bar{B}_{c,2}}{1 + B_c}.
\end{equation}
This target rewards cells with low local branch count but high surrounding
banking density. A gradient boosting regressor is trained to predict
$O_c$, and predictions are normalized to a 0--100 opportunity score.

\subsection{Interpretability of the Opportunity Score}
The opportunity score is not intended to be a causal estimate of future
branch profitability. It is an interpretable prioritization function. A
high-scoring cell has three practical meanings: local branch supply is low,
nearby market activity is nontrivial, and the model predicts a strong
spatial mismatch between local and surrounding supply. This makes the score
useful for analyst triage, site-screening workflows, and visualization.

\section{System Architecture}
GeoRisk-IQ is implemented as a modular Python pipeline with four layers:
data ingestion, spatial feature generation, machine-learning scoring, and
interactive visualization. The ingestion layer reads FDIC files from the
project root and normalizes records into pandas data frames. The spatial
layer maps coordinates to cells, computes local and neighboring density,
and aggregates financial indicators. The modeling layer trains regression
and classification models. The visualization layer exposes results through
a Streamlit dashboard using Plotly and pydeck.

\begin{figure}[!htbp]
\centering
\figimage{\columnwidth}{fig5_architecture.jpg}
\caption{GeoRisk-IQ system architecture. FDIC records are cleaned,
assigned to spatial cells, transformed into neighborhood features, and
used by scoring and validation models.}
\label{fig:architecture}
\end{figure}

\subsection{Operational Workflow}
In operational use, the pipeline can be rerun whenever updated FDIC records
are downloaded. Analysts can inspect the top opportunity table, open the
opportunity map, compare the full density surface, and run the IEEE
validation module when failure and financial files are present. This design
keeps the research model and the dashboard aligned: the same features used
for training are shown in the interface.

\begin{table}[t]
\caption{Pipeline Stages}
\label{tab:pipeline}
\centering
\begin{tabular}{lll}
\toprule
Stage & Input & Output \\
\midrule
Ingestion & FDIC files & Clean data frames \\
Grid build & Coordinates & Spatial cells \\
Feature build & Cells & Neighborhood matrix \\
Scoring & Feature matrix & Opportunity ranks \\
Validation & Labels + features & Model metrics \\
Dashboard & Results & Maps and tables \\
\bottomrule
\end{tabular}
\end{table}

\section{Methodology}
\subsection{Failure Labels}
For risk validation, each cell receives a binary label:
\begin{equation}
    y_c =
    \begin{cases}
    1, & \text{if at least one mapped bank failure occurs in cell } c,\\
    0, & \text{otherwise.}
    \end{cases}
\end{equation}

\subsection{Model Families}
We compare four model families:
\begin{enumerate}
    \item Logistic regression using the full feature set and balanced class
    weights.
    \item Spatial-only gradient boosting using neighborhood and density
    features.
    \item Financial-only gradient boosting using cell-level financial
    aggregates.
    \item Hybrid gradient boosting using both spatial and financial
    features.
\end{enumerate}
All experiments use a fixed random seed of 42.

\begin{table}[t]
\caption{Model Configuration}
\label{tab:model_config}
\centering
\begin{tabular}{ll}
\toprule
Component & Setting \\
\midrule
Opportunity model & GradientBoostingRegressor \\
Validation model & GradientBoostingClassifier \\
Cross-validation & 10-fold stratified \\
Random seed & 42 \\
Scaling & StandardScaler before modeling \\
Imbalance handling & Balanced sample weights \\
Permutation test & 10,000 swaps \\
Bootstrap CI & 1,000 resamples \\
\bottomrule
\end{tabular}
\end{table}

\subsection{Evaluation Protocol}
The opportunity scorer is evaluated on a held-out test split using mean
absolute error and $R^2$. The failure-risk task is evaluated with 10-fold
stratified cross-validation, temporal holdout, bootstrap confidence
intervals, permutation tests, and ablation. For temporal validation, the
model trains on pre-2010 failures plus clean cells and tests on post-2010
failures plus held-out clean cells. The classification threshold is chosen
to maximize minority-class F1 on the validation loop.

\subsection{Statistical Testing}
The permutation test compares two models by randomly swapping predicted
scores sample-by-sample and recomputing the AUC difference. This directly
tests whether the observed difference between models is larger than would
be expected if their predictions were exchangeable. Bootstrap confidence
intervals are computed by resampling concatenated cross-validation
predictions while preserving the observed score-label relationship.

\section{Results}
\subsection{Opportunity Scoring}
Table~\ref{tab:opportunity_metrics} reports opportunity-scoring results.
The high $R^2$ should be interpreted as fit to a deterministic proxy target,
not as proof of observed future branch openings. The value of this module is
to produce a transparent ranking of underserved cells near dense banking
activity.

The highest-ranked opportunity cells are not necessarily the largest
markets. Instead, they are locations with a favorable contrast between
their own branch count and surrounding density. This property is important
for banking strategy because it surfaces candidate gaps rather than simply
rediscovering already dense markets.

\begin{table}[t]
\caption{Opportunity Scoring Performance}
\label{tab:opportunity_metrics}
\centering
\begin{tabular}{lc}
\toprule
Metric & Value \\
\midrule
Institutions loaded & 27,832 \\
Spatial cells & 6,086 \\
Mean absolute error & 0.0081 \\
$R^2$ & 0.9991 \\
\bottomrule
\end{tabular}
\end{table}

\begin{table}[t]
\caption{Example Top Opportunity Cells}
\label{tab:top_cells}
\centering
\begin{tabular}{lccc}
\toprule
Cell center & Score & Branches & Assets (\$K) \\
\midrule
(41.4, -88.0) & 100.0 & 1 & 260,568 \\
(29.4, -95.4) & 84.4 & 1 & 12,268 \\
(32.6, -96.6) & 82.9 & 1 & 376,443 \\
(30.0, -95.8) & 75.8 & 1 & 27,025 \\
(39.2, -77.0) & 75.4 & 1 & 13,765,535 \\
\bottomrule
\end{tabular}
\end{table}

\begin{figure}[!htbp]
\centering
\figimage{\columnwidth}{fig1_hex_map.jpg}
\caption{FDIC branch density and opportunity-score dashboard. The density
view uses spatial cells to summarize local branch counts and surrounding
opportunity signals.}
\label{fig:hex_map}
\end{figure}

\begin{figure}[!htbp]
\centering
\figimage{\columnwidth}{fig2_top_opportunities.jpg}
\caption{Top opportunity zones. Filled red cells mark high-ranked candidate
areas with low local supply and high surrounding banking density.}
\label{fig:top_opportunities}
\end{figure}

\subsection{Bank-Failure Cross-Validation}
Table~\ref{tab:cv_results} compares the four model families. The hybrid
model achieves the highest ROC-AUC. Logistic regression achieves the
highest F1 in this configuration, indicating that thresholding and class
imbalance remain important practical considerations.

\begin{table}[t]
\caption{10-Fold Cross-Validation on Failure Labels}
\label{tab:cv_results}
\centering
\begin{tabular}{lcc}
\toprule
Model & ROC-AUC & F1 \\
\midrule
Logistic regression & 0.8143 & 0.5780 \\
Spatial-only GBM & 0.8279 & 0.5354 \\
Financial-only GBM & 0.6029 & 0.1697 \\
\textbf{Hybrid GBM} & \textbf{0.8341} & 0.5320 \\
\midrule
\multicolumn{3}{l}{Hybrid AUC 95\% CI: [0.8226, 0.8454]} \\
\bottomrule
\end{tabular}
\end{table}

The hybrid model improves over the spatial-only model by
$\Delta$AUC = +0.0062. Although this improvement is numerically modest, a
two-sided permutation test on concatenated cross-validation predictions
finds it statistically significant ($p=0.0013$). The hybrid model also
substantially improves over the financial-only model
($\Delta$AUC = +0.2312, $p<0.0001$). A supplementary Wilcoxon signed-rank
test on fold-level AUCs gives $p=0.0020$ for hybrid versus spatial-only.

The model comparison yields two observations. First, spatial-only features
are strong enough to produce competitive AUC, confirming that neighborhood
structure has predictive signal. Second, the hybrid model provides the best
ranking quality, showing that financial and spatial features are
complementary. The financial-only model is weaker in this run, which is
consistent with the snapshot-based financial data limitation described in
the data section.

\begin{figure}[!htbp]
\centering
\figimage{\columnwidth}{fig2_roc_pr.jpg}
\caption{ROC and precision-recall curves for the four banking-domain
models.}
\label{fig:roc_pr}
\end{figure}

\subsection{Temporal Holdout}
Table~\ref{tab:temporal} reports temporal holdout performance. The model
trains on pre-2010 failures and clean cells, then tests on post-2010
failures and held-out clean cells.

Temporal validation is stricter than random cross-validation because it
asks the model to generalize from earlier failures to later failures. The
holdout AUC of 0.7331 is lower than the cross-validation AUC, as expected,
but remains above random ranking and demonstrates that the learned signal
does not depend entirely on random fold mixing.

\begin{table}[t]
\caption{Temporal Holdout Results}
\label{tab:temporal}
\centering
\begin{tabular}{lc}
\toprule
Metric & Value \\
\midrule
Precision & 0.3776 \\
Recall & 0.4512 \\
F1 & 0.4111 \\
ROC-AUC & 0.7331 \\
\bottomrule
\end{tabular}
\end{table}

\subsection{Feature Attribution}
SHAP attribution for the hybrid model identifies the second-ring density
ratio as the strongest feature, followed by first- and second-ring neighbor
density and mean financial assets. This finding supports the core
hypothesis that medium-range spatial market context contributes information
that is not captured by local counts alone.

\begin{table}[t]
\caption{Top Hybrid Model Feature Attributions}
\label{tab:importance}
\centering
\begin{tabular}{lr}
\toprule
Feature & Mean absolute SHAP \\
\midrule
density\_ratio\_k2 & 1.0005 \\
neighbor\_avg\_k1 & 0.3804 \\
neighbor\_avg\_k2 & 0.3253 \\
fin\_asset\_mean & 0.3103 \\
density\_ratio\_k1 & 0.2706 \\
asset\_total & 0.2604 \\
neighbor\_avg\_k3 & 0.1316 \\
is\_isolated & 0.0711 \\
fin\_roa\_mean & 0.0486 \\
fin\_netinc\_mean & 0.0360 \\
\bottomrule
\end{tabular}
\end{table}

\begin{figure}[!htbp]
\centering
\figimage{\columnwidth}{fig3_feature_importance.jpg}
\caption{Feature attribution for the hybrid model. Spatial density ratios
and neighborhood averages are prominent contributors.}
\label{fig:feature_importance}
\end{figure}

\subsection{Ablation}
Table~\ref{tab:ablation} summarizes a 5-fold ablation study. Removing all
neighbor features produces the largest F1 drop, supporting the role of
spatial context in failure-risk validation.

\begin{table}[t]
\caption{Ablation Study, 5-Fold CV F1}
\label{tab:ablation}
\centering
\begin{tabular}{lcc}
\toprule
Configuration & F1 & $\Delta$F1 \\
\midrule
Full model & 0.5037 & 0.0000 \\
Without $k=1$ ring & 0.5040 & +0.0003 \\
Without $k=2$ ring & 0.4887 & -0.0150 \\
Without $k=3$ ring & 0.5059 & +0.0022 \\
Without all neighbors & 0.3564 & -0.1473 \\
Without financial & 0.4957 & -0.0080 \\
\bottomrule
\end{tabular}
\end{table}

\begin{figure}[!htbp]
\centering
\figimage{\columnwidth}{fig4_hrsa_map.jpg}
\caption{Ablation dashboard view. Removing all neighbor features produces
the largest F1 reduction, supporting the role of spatial context in the
hybrid validation model.}
\label{fig:ablation}
\end{figure}

\subsection{Runtime and Practicality}
The banking opportunity pipeline completes quickly enough for interactive
use because it trains on thousands of spatial cells rather than millions of
raw records. The full validation pipeline is heavier because it includes
cross-validation, bootstrap intervals, permutation testing, SHAP
attribution, and ablation. In practice, opportunity scoring can be run
frequently, while the full validation module can be scheduled as a monthly
or quarterly audit.

\section{Discussion}
The results show that spatial neighborhood features are valuable for both
opportunity scoring and failure-risk validation. The opportunity scorer
surfaces cells where local branch supply is low relative to surrounding
activity, while the risk-validation model shows that spatial context
improves ranking quality over a financial-only baseline.

The strongest attribution result is the second-ring density ratio. This
suggests that medium-range market context may better capture competition
and isolation than immediate local branch counts. In practical terms, a
branch or market cell should not be evaluated only by its own local density;
its relationship to surrounding cells is also informative.

The financial-only model is weaker than expected in the current
reproducible run. This likely reflects the use of a latest-quarter financial
snapshot instead of full historical lagged financial trajectories. Future
work should incorporate time-aware financial features, rolling trends, and
survival modeling.

\subsection{Why Spatial Ratios Matter}
The second-ring density ratio is the strongest attribution feature in the
current hybrid model. This is intuitive in branch-network analysis: the
immediate cell captures local supply, but the surrounding ring captures the
market context in which customers, competitors, and service gaps operate.
A branch in a low-density cell near high-density neighbors may indicate an
underserved boundary market. A branch in a low-density cell surrounded by
other low-density cells may instead indicate rural sparsity. The ratio
distinguishes these cases.

\subsection{Dashboard Implications}
The dashboard is not just a presentation layer. It provides a feedback loop
for validating whether model-ranked cells make geographic sense. Analysts
can compare the opportunity map with the density map, inspect raw cell
features, and download scored grids. This visual inspection is especially
important for geospatial models, where coordinate errors or boundary
artifacts can create misleading rankings.

\section{Deployment Considerations}
GeoRisk-IQ can be deployed as a batch scoring tool or as an analyst-facing
dashboard. In batch mode, updated FDIC files are downloaded, cells are
rebuilt, models are retrained, and ranked opportunity zones are exported.
In dashboard mode, users can alter grid resolution and top-$N$ thresholds,
then inspect how the ranking changes. A conservative deployment should
include checks for missing files, invalid coordinates, class imbalance, and
large shifts in feature distributions across runs.

\begin{table}[t]
\caption{Operational Checks Before Deployment}
\label{tab:deployment_checks}
\centering
\begin{tabular}{ll}
\toprule
Check & Reason \\
\midrule
Coordinate validity & Prevent false cells \\
Record-count drift & Detect incomplete downloads \\
Failure prevalence & Monitor label stability \\
Financial snapshot date & Avoid stale indicators \\
Top-cell review & Catch geographic anomalies \\
Metric regression & Detect model degradation \\
\bottomrule
\end{tabular}
\end{table}

\section{Limitations}
This study has several limitations. First, the current implementation uses
an H3-style deterministic grid abstraction rather than a production native
H3 dependency. Second, failure locations are inferred through available FDIC
certificate and coordinate mappings; records without reliable geospatial
links are omitted. Third, opportunity scoring uses a proxy target rather
than observed future branch openings. Fourth, the financial feature set is
cell-level and snapshot-based. Finally, figures included in this manuscript
should be regenerated from the exact experimental run used for submission.

There are also evaluation limitations. The opportunity target is synthetic,
so its performance metrics measure fidelity to a scoring rule rather than
observed business outcomes. The failure-risk labels are historical and may
reflect regulatory, macroeconomic, and institutional processes not captured
by the current feature set. The model should therefore be interpreted as a
decision-support tool, not as an autonomous approval or supervisory system.

\section{Future Work}
Future work will focus on four extensions. First, replacing the H3-style
grid abstraction with native H3 cells will improve geometric fidelity and
allow easier comparison with other geospatial systems. Second, historical
financial trajectories should be incorporated through lagged features,
rolling means, and survival-model variants. Third, opportunity scoring
should be validated against observed branch openings, closures, deposits,
or market-share changes. Fourth, the cross-domain healthcare pipeline can
be restored when the HRSA source file is available and evaluated with the
same statistical discipline used in the banking experiment.

\section{Conclusion}
GeoRisk-IQ provides a reproducible spatial-financial framework for bank
branch opportunity scoring and failure-risk validation. The system converts
FDIC institution records into spatial cells, computes neighborhood density
features, and combines them with financial aggregates. In the current
reproducible banking experiment, the hybrid model achieves ROC-AUC of
0.8341, with statistically significant improvement over spatial-only and
financial-only baselines. These results support further development of
geospatial financial intelligence systems for branch network planning,
supervisory analytics, and market-access analysis.

\section*{Acknowledgment}
The author thanks the public data providers whose open datasets make
reproducible geospatial financial analysis possible.

\appendix
\section{Reproducibility Notes}
The current implementation uses Python, pandas, NumPy, scikit-learn,
Plotly, pydeck, Streamlit, SciPy, and SHAP. Random seeds are fixed with
\texttt{RANDOM\_STATE=42}. FDIC institution, failure, and financial files
are loaded from the project root. The Streamlit dashboard can be run with
\texttt{streamlit run app.py}.

\begin{thebibliography}{10}

\bibitem{fdic}
Federal Deposit Insurance Corporation, ``BankFind Suite API.''
[Online]. Available: \url{https://banks.data.fdic.gov/}

\bibitem{h3}
Uber Technologies, ``H3: A Hexagonal Hierarchical Geospatial Indexing
System.'' [Online]. Available: \url{https://h3geo.org/}

\bibitem{cole1998}
R. A. Cole and J. W. Gunther, ``Predicting bank failures: A comparison
of on- and off-site monitoring systems,'' \textit{Journal of Financial
Services Research}, vol. 13, no. 2, pp. 103--117, 1998.

\bibitem{betz2014}
F. Betz, S. Oprica, T. A. Peltonen, and P. Sarlin, ``Predicting distress
in European banks,'' \textit{Journal of Banking and Finance}, vol. 45,
pp. 225--241, 2014.

\bibitem{tian2015}
S. Tian, Y. Yu, and H. Guo, ``Variable selection and corporate bankruptcy
forecasts,'' \textit{Journal of Banking and Finance}, vol. 52, pp. 89--100,
2015.

\bibitem{friedman}
J. H. Friedman, ``Greedy function approximation: A gradient boosting
machine,'' \textit{Annals of Statistics}, vol. 29, no. 5, pp. 1189--1232,
2001.

\bibitem{shap}
S. M. Lundberg and S.-I. Lee, ``A unified approach to interpreting model
predictions,'' in \textit{Advances in Neural Information Processing
Systems}, 2017, pp. 4765--4774.

\bibitem{sklearn}
F. Pedregosa et al., ``Scikit-learn: Machine learning in Python,''
\textit{Journal of Machine Learning Research}, vol. 12, pp. 2825--2830,
2011.

\end{thebibliography}

\end{document}
